\subsubsection{C 类：数据口径、隔离验证与决策解释并重}

C 类从字段、样本单位、时间范围、缺失异常和训练/验证边界写起。有效表达包括
“以……为一条观测，按主体/时间划分训练与验证集”“所有标准化和特征选择仅在
训练折拟合”“相对持续性/线性基线，模型在……指标上……”“该变化主要对应……
特征，但结论限于……范围”。评价排序需解释指标正向化、权重来源和排序稳定性；
预测需把误差转换为补货、时点或风险决策的影响。

需要避开的是全数据清洗后再切分、用随机切分检验未来预测、把相关性写成因果、
只报告最优模型以及用二维降维图证明分类性能。图表以数据质量、分布与相关、
验证协议、残差/校准、特征影响和决策结果为主。

C006、C085、C169、C283、C155、C229、C050、C126、C228、C109、C142、C170、C227、
C305、C038、C063、C094、C234、C023 与 C132 共二十篇 C 题均已通过人工检查。以下归纳
同时覆盖业务统计、规划决策和医学筛查，但每种句式仍须回到本题的样本单位、时间口径和
决策对象，不能把不同应用中的模型名机械拼接。
C006 的正向价值首先体现在“把业务语言改写成可计算量”的过程。论文没有直接把“供货
连续性”当成一个无来由的分数，而是拆成间隔次数、平均间隔周数和平均连续供货周数，
再说明各指标的成本/效益方向、归一化和熵权。这里可迁移的不是熵权--TOPSIS 这个名称，
而是“业务概念--不可直接观测--代理指标--方向--权重--解释”的完整段落。

该篇在周供货能力处还主动拒绝了省事的历史均值：均值会把不同周次抹平，经济作物的
产量又具有周期性，故改为在十个 24 周周期中取相同周次的历史位置，再用 GP 估计订货量
对应的供货误差。逐周订购时，上周库存、实际供货和消耗量共同形成下一周初始状态；
二元供应商筛选与逐周订购使用遗传算法，而转运环节按损耗率、整批供应、剩余运力和 A 类
优先等规则分配。这样的写法把三个求解角色分开，避免只因同属一条业务链便统称“用 GA
优化”。

结果解释也保留了人的判断痕迹。作者先给 24 周总体方案，再选第 3 周和第 13 周说明库存、
供货与转运怎样在具体周次闭合；发现原方案没有表现出材料偏好后，先诊断成本系数所对应的
替代机制，再只调整 A/C 原料的等产能约束。调整后的成本最大约降 $5\%$、平均仅降
$0.956\%$，论文承认改善幅度有限，并把原因落到供应商数量少、供货能力贴近需求和可行域
狭窄，而不是把正向变化包装成“显著提升”。问题四沿用转运模型时，又说明明输入、输出和
优化目标均未改变；这比一句“沿用上问模型”更能说明模块为何可复用。

C085 把另一类业务判断写得具体。四个供应商指标不是只给名称和权重，而是逐项解释
“统计量怎样定义、数值方向表示什么、它对应哪一种供货能力”；稳定供货量又先受两条
规则筛选：供货/订货超过 $0.9$，且供货量不超过同周期位置均值的三倍，只有同时满足时
才取最大候选值。这样的段落先让阈值承担稳定性与异常排除作用，再给期望量，不以
“经预处理得到”省略判断过程。

算法改进也由可指认的失败触发。固定周权重可能使总体供货看似充足，却让个别周严重
缺货，所以作者按现有供货与期望产能的缺口更新下一轮权重；第 7、8 周在数据中连续处于
淡季，因而前视窗口取两周并提前备料。转运部分则先拆分超过单个转运商容量的大项，再按
可靠性和材料价值做 0--1 背包，最后对未装入项贪心补救。结果比较时，历史采购成本虽低，
但历史平均产能没有达到 28200，故改用采购成本/产能比较 $0.722032258$ 与 $0.721501175$。
这几处都先写“原口径或原算法具体在哪儿失真”，再给修正动作；比“综合考虑多种因素，
采用混合算法优化方案”贴近人类作者的真实推理。

C169 展示了另一种容易被 AI 文本抹掉的动作：作者会说明哪些观察\emph{不进入}模型。
前 50 家供应商的序列中既有周期性异常，也有非周期尖峰；论文先解释两类现象的可能成因，
随后用“尖峰会干扰周期规律”和“新订货方案会打破旧规律”两个理由，明确下文不另行建模。
类似地，初始库存没有方便地设为 0，而是以“企业并非刚刚成立”为起点，回到历史 24 周
周期的同位状态估计。转运模型也先判断损耗无记忆，确认过去损耗不影响本周后，才按周拆分；
这种次序比“为简化计算，将问题分解为 24 个子问题”更可信。

该篇还把业务后果直接塑造成数学结构。缺供对生产的危害重于超供，作者先列分段函数应满足
的符号、曲率和两侧大小条件，再构造 $e^x-1$ 与 $e^{-2x}-1$，而不是先选函数再补解释。
候选集的前 50 家又由“覆盖总供货量 90\% 以上、尾部仅个位数、继续增加会损害经济性”共同
限定。跨周聚合时，最少供应商取 24 个逐周最小数的最大值，最大产能取 24 个逐周上界的
最小值；聚合方向不同，是因为前者要找统一可用的供应商数，后者要保证每周都能达到产能。

C283 把评价指标进一步写成附件上的实际操作。忙期均量并非对 240 周直接求平均，而是先去掉
无订单形成的零周；信誉度又让相对订供偏差乘以订货量的对数，使大额且按量履约的订单提供更多
证据。熵权结果出现供给弹性等业务重要指标权重过低的情形后，作者没有用一句“主客观结合”
带过，而是先将指标归到绩效、战略意义和潜力，再在适度范围内用 AHP 修正并通过一致性门。
这形成“数据权重--具体异常--业务冲突--有限修正--一致性检查”的完整段落。

双层贪心的选型和停止也保留了真实取舍。论文先写线性规划的变量/约束关系，随后指出动态规划的
状态空间约为 $6000\times50\times24$，随后利用供应商和转运商已有排序构造两层贪心。库存达到
两周需求后立即停止，不是因为程序“收敛”，而是继续下单会转向较差供应方或运输方并增加存运
成本。四个关键订货式共用历史忙期均量，再按约束、成本、初期稳健性和现实含义逐个淘汰；随机
运行则明确只取“频率较高且目标较优”的典型代表，不把一次输出包装成唯一最优。

问题三的类别权重同时写了上下边界：过小不能突出 A/C 的存储差异，过大又会牺牲采购经济性和
合作可靠性。到问题四，目标由成本/损耗最小转为产能最大，先前排名靠后的供应商也可能贡献供货，
所以候选从前 50 家扩到全部 176 家。结果段发现实际停止条件与预期瓶颈不一致时，作者保留
疑问并给出“提高进货参数”或“降低偏差率”两种解释；生产曲线与接收曲线同向、约滞后 3 周，
则把库存传导写成了可读的时间关系。

C155 先处理“缺失到底表示什么”，而不是把所有空格交给同一种插补。化学成分空白按题意表示
未检出，故填为 0；颜色字段却不能照此处理，因为相同已知条件仍可能对应多种颜色，且表面随机
取样不足以代表整件文物。论文据此拒绝猜测性填补，只在涉及颜色的分析中忽略 19、40、48、
58 号样品。类似地，成分总和低于 $85\%$ 的 15、17 号点先由题面阈值判为无效，再剔除；其余
样本闭合到 $100\%$ 后才作中心化对数比变换。这种“题面语义--具体编号--判断--动作--影响范围”
比一句“对缺失值和异常值进行预处理”更能使读者知道数据发生了什么。

统计方法也由前提而不是方法名分流。作者先列联表和期望计数，满足 Pearson 卡方检验前提的组
直接使用 Pearson，不满足的组改用 Yates 校正；风化与类型的 $p=0.009$，与颜色、纹饰的
$p$ 分别为 $0.507$、$0.084$。随后 $k=2$ 的 Q 型聚类借助前问趋势和已知标签命名风化/未风化
簇，又用聚类结果反向核对前问判断；$k=4$ 随后将风化展开为未风化、轻度、中度和重度四阶段。
同一条证据在前后问题之间双向使用，使各问不是互不相干的算法堆叠。

预测段先逐成分比较二次回归的 $R^2$，只预测拟合较好的成分，对较低者采取不变策略；若异常点
使“曲线到观测值的距离”被放大，则只对这个唯一敏感量作单调压缩，再经逆对数比和闭合约束恢复
总和为 $100\%$ 的成分向量。分类段又把 R 型和 Q 型聚类的职责拆开：R 型先将十四个相关成分
变量分组并选代表量，Q 型再用代表量划分样品亚类。未知样品由两棵按风化状态分层的 PbO 决策树
分类，并以 Q 聚类交叉核对；最后固定每一种成分为灰色关联母序列，再比较两类玻璃的最大、最小
关联及分布集中程度。这里可迁移的是“指标门控--局部修正--闭合验证”“变量先降相关、样本后分组”
以及“同一基准下比较极值和分布”，不是把所有方法名称并列到摘要中。

C229 进一步说明，即使缺失的是同一个颜色字段，填补方法也不必一刀切。论文先用纹饰、表面
风化与玻璃类型组成匹配条件，再逐个核对候选样品的参考质量。40、58 号分别能找到成分相近且
可参照的 39、11 号，因此用热卡填为深绿、浅蓝；19、48 号所在组虽有匹配对象，但多数化学
样本来自未风化采样点，参考性较低，故转用众数并填为浅蓝。正文把“分组条件--候选编号--
参考质量--方法分支--填补结果”连续写出，比笼统声称“对定性变量采用合理插补”更可信。

化学成分的空白与数值 0 也先按数据生成过程解释：前者来自低于检测限，后者可能是保留位数后的
舍入结果，二者都不能当作普通零值进入对数比。作者以列最小正值近似检测限，取
$\delta_j=0.65e_j$ 作乘法替换，随后将向量闭合到 $100\%$。SiO2 占比过大而遮蔽其他成分时，
论文保留全部成分图，又补一张去除 SiO2 的图，最后写“结合两图进行分析”。这种做法没有为了
看清小量而删掉主量证据，也没有让一张比例失衡的图承担全部结论。

分类段先用决策树给出可复核核的 PbO=1.5 阈值，在 47 个训练样品和 20 个测试样品上报告
测试准确率 $100\%$；随后 PLS-DA 以 PbO、K2O、BaO、SrO 的 VIP 值复核多变量差异。亚类
没有已知标签，作者据此明确把任务判为无监督学习，再由 WSS 曲线选 $k=3$。铅钡类三簇界限清楚，
高钾类没有明显分界，正文只写“仅可粗略将其分为三类”。这里可迁移的是“可解释规则--多变量
复核”和“无标签身份--诚实限定聚类强度”，而不是把两种模型或两类图统一包装成效果良好。

未知样品的 PLS 维数同时比较十折随机段交叉验证误差与累计解释率：三个成分时 CV RMSEP 为
$0.1371$，继续增加到四个成分只到 $0.1372$，而 X 与类别累计解释率已达 $93.55\%$ 与
$92.00\%$，故停止在三维。A1、A6、A7 判为高钾，其余五件判为铅钡，树与 PLS 结论一致。
问题四又从数据结构推出检验：相关矩阵对称，故只取上三角；相同成分对一一对应，故用配对
Wilcoxon，$p=0.905$。参数选择和检验选择都由可检查的两层理由推进，而非由方法名称代替判断。

C050 把数据侧写写成模型入口，而不是把探索性图表集中摆放后便不再使用。论文先从品种丰富度、
季度销量、损耗、描述统计和时序变化观察附件，并在负价记录的业务含义明确后剔除退货数据；后续
相关分析、半年预测窗口、可售单品范围和优化约束都能回指这些观察。单品数量过多时，作者也不是
随手挑几条曲线，而是先规定各类展示总销量最高者，以及总销量最低但不少于 $2000\,\mathrm{kg}$
者，其余结果转入附件。这样的顺序形成“侧写对象--实际观察--下游判断”和“代表规则--正文展示--
全量附件”的双重接口，使图表选择本身可以复核。

回归和人工特征的解释继续落回业务口径。论文没有停在相关系数、显著性或线性方程，而是逐类说明
销量每增加 $1\,\mathrm{kg}$ 时成本加成定价相应增加或减少多少元。问题四又分别以品类单品种数、
季度中心附近的节气和主要节日构造丰富度、季节与节日指标，并先说明代表对象、计算口径及其所代理
的经营现象，再进入灰色关联。可迁移的是“系数--业务单位--方向--决策含义”和“对象--口径--代理
语义--下游用途”，不是在结果段重复一遍公式或把所有人工指标笼统称为特征工程。

该篇还把现有附件能够支持的结论与尚待采集的信息分开。附件内的销量、价格和日期先完成三项关联
分析；客流、竞品、反馈、天气、供应链等信息则另设段落，按“如何采集--能够区分什么机制--支持
何种经营动作”展开。打折频率较低既可能表示商品受欢迎或供不应求，也可能只是促销必要性较弱；
频率较高则可能来自销售乏力或保鲜期较短，故需联查销量、反馈和竞争情况后再决定定价或促销。
这里的关键不是列更多外部变量，而是保留“代理量--竞争机制--追加证据--分情形动作”的推理过程。
正文只保留代表图表、关键数值、选择规则和解释，完整单品表、补充图与代码进入附录，使证据分流与
版面分工保持一致。

C126 的第一条写作价值，是把方法切换前的压力写全。日销售量不服从正态、零值多、三年序列有
季节结构，OpenBUGS 的 MCMC 又有实际计算负担；作者因此不是直接宣布“采用 Bayesian 方法”，
而是先用同期平均把三年数据压成一年，再交代新口径怎样进入先验、迭代、燃烧期和后验诊断。
可迁移的结构是“原数据现象--旧口径的具体困难--转换规则--新输入--求解设置”，它允许读者看见
为什么换模型，而不是用“提高精度、降低复杂度”两个空泛评价代替选择过程。

第二条价值来自业务事实对模型结构的授权。论文先写蔬菜当日未售后隔日大多难以原价销售，由此
把有效记忆限定为两天，再选择二阶 VAR；模型阶数不是孤立的调参结果。预测出现负数时，作者也没有
机械删去，而把它解释为未售且未退的余货、次日可折价以及当日无需补货。这样，“业务事实--时间
尺度--模型阶数”和“数学读数--库存状态--经营动作”各自形成完整自然段，比“综合考虑实际情况”
更能留下人的判断痕迹。

第三条价值是信息压缩和证据分流。相关系数、阈值与显著性先被压成 T/F 判断矩阵；13 维或更大的
结果难以放入正文时，作者明确以花叶类或 7 个水生根茎类为例，正文保留判断规则、代表参数表、
MCMC 诊断和模型评价，完整图表、四组 BUGS 数据块、EViews 工程、R 检验与两问 LINGO 代码进入
支撑材料。VAR 路线筛得 29 个单品后，又从半小时销售频率独立筛选并得到近似名单；这一动作只称
有限交叉核对。预测量随后进入利润收入项，成本、损耗和利润率进入约束，问题三随后将最近 7 日反推
进货量、平均计划量、前日余货、新补货和 $2.5\,\mathrm{kg}$ 陈列下限逐项接起。模型之间因此有
可见中间量，而不是几个算法名称排成一条箭头。

C228 把开放任务的第一步写成“先决定看什么”。题目没有指定具体探究方向，作者没有用“从多个
维度深入分析”一笔带过，而是把品类和单品分别放到时间分布、销售分布及相互关系三个观察面上，
并说明这些角度合在一起怎样覆盖销售规律。年、周、日周期的统计读数之后又接上采收季、春节、
周末采购和上下班节奏；图形不作孤立展示，随后用于确定预测窗口，并解释相应的经营现象。

同篇还保留了两次不同性质的方法退让。第一次是试过 ARIMA、回归和灰色预测后，写明当前
序列上的效果不理想，再改用 LSTM；第二次是面对近七日单品的缺失、短时无明显趋势和随机波动，
反而不再强行使用复杂预测，而以加权平均给出当前时窗内的估计。前者是“候选--表现--换模”，
后者是“数据能力--复杂方法不再合适--主动降级--限定时窗”，两种取舍都比事后罗列算法优点贴近
真实建模过程。

预测到优化的接口也没有被一句“代入模型”吞掉。LSTM 先给基线销量，双对数弹性再按价格变化
修正需求，修正量进入收入项，成本、损耗和补货量分别进入目标或约束；问题三新增是否进货的
$0/1$ 变量、连续加成率和补货量以及双目标，作者因变量结构改变才从连续 PSO 转向混合编码的
NSGA-II。结果段又把方案放回同季历史日利润区间比较。最后的补采建议从成本导向定价的边界长出，
具体指向竞品价格、消费者预期和收入，而不是另列一份与前文无关的“更多影响因素”。

C109、C142、C170、C227 与 C305 都把企业画像、风险估计和信贷方案串在一条决策链上。
这组论文最值得保留的不是某个分类器名称，而是分问之间的接口意识：发票先整理成企业级特征，
特征产生风险或等级，风险再进入准入、额度、利率和冲击重算。写作时应逐项说明“上一问的哪个
输出成为本问的哪个输入”，并让长表和附录代码承担逐企业追踪，正文只保留规则、代表结果和解释。

这五篇的自然行文还体现在业务量会被重新翻译。交易连续性不是抽象分数，而具体到交易次数、
金额、波动或中断；外部冲击不是一句“考虑疫情影响”，而具体到行业、地区、需求或偿付能力的
参数变化。段落宜写成“由附件可计算……，故以……刻画……；该量进入……项，从而影响……”。
如果只能给单一情景，就把结论限定在该情景，不用“具有普遍适用性”之类空泛结尾。

C038、C063、C094 与 C234 展示了种植规划中四种互补的表达。C038 把七年规划先写成明确的
地块、季次、作物和轮作约束，再用 CVaR 讨论尾部风险；C063 让典型情景负责展示候选差异，
却把最终裁决放到共同随机环境中；C094 将“优先利用、稳定种植、兼顾收益”翻译为管理参数
$p/q$；C234 先从数据表确定索引，随后将三类真实邻接关系写进染色体和修复算子。共同环境复评、
管理语言参数化和可行编码，都是“现实判断怎样进入模型”的具体答案。

这组规划论文的结果段也不把“收益更高”单独拿出来。作者会同时写期望收益、最坏情形或尾部
风险，说明为了稳健性牺牲了多少利润；候选差别很小时，回到活跃约束、种植集中度和实际执行
难度裁决。适合复用的句序是：“方案……在基准情景下……；置于相同的……情景后，……；
差异主要来自……约束，因此选择……。”三个省略处都必须填入本题对象和读数。

C023 把医学筛查中的统计模型与行动后果分开。文章先用混合效应模型处理重复测量，随后将达标概率
送入分组和检测时点决策；安全、及时与可执行分别对应不同损失，概率拟合好并不自动等于决策好。
“结果可疑”也被翻译为复测或升级处置。这里可照搬的组织是“概率估计--阈值含义--错误代价
--行动规则--稳健核对”，而不是把 Logit、动态规划和 FDR 并排罗列。

C132 先将多次检测记录从“行”还原为“孕妇”，再定义最早达标时点，这是其全文最关键的口径
转换。缺测时按实测证据、保守估计和无法判断三级回退；首次就诊即可获得的静态特征交给 MLP，
随孕周变化的序列交给 LSTM，存在右删失的时点风险交给 DeepHit。作者没有用“多模型融合”
概括这些角色，而是逐一说明各模型接住哪类信息和哪种缺失。

该篇对结果突变的解释也值得保留：当达标概率接近 $90\%$ 阈值时，等待的边际代价突然增大，
于是曲线变陡；最终融合不是把各专家在全数据上的输出直接重用，而以折外概率作为元分类器输入。
结果段同时报告总体表现、关键少数类召回和对应样本数，使“整体不错”不会遮住真正需要识别的
高风险人群。
